\section{Gambaran Umum Sistem}

Aplikasi yang dikembangkan merupakan sistem berbasis web yang bertujuan untuk melakukan pencarian elemen HTML pada struktur DOM berdasarkan CSS selector. Sistem ini memanfaatkan algoritma Breadth First Search (BFS) dan Depth First Search (DFS) sebagai metode traversal.

\begin{tcolorbox}[colback=gray!10,title=Input dan Output Sistem]
\textbf{Input:}
\begin{itemize}
    \item URL website atau teks HTML
    \item CSS Selector
    \item Pilihan algoritma (BFS / DFS)
    \item Jumlah hasil
\end{itemize}

\textbf{Output:}
\begin{itemize}
    \item Elemen hasil pencarian
    \item Visualisasi DOM
    \item Jalur traversal
    \item Waktu eksekusi
\end{itemize}
\end{tcolorbox}

% --------ARSITEKTUR SISTEM--------

\section{Arsitektur Sistem}

Sistem dibangun dengan pendekatan client-server.

\begin{center}
\begin{tabular}{|c|c|c|}
\hline
Frontend & $\rightarrow$ & Backend \\
\hline
Input user & & Parsing + BFS/DFS \\
\hline
Visualisasi & $\leftarrow$ & Hasil traversal \\
\hline
\end{tabular}
\end{center}

\begin{itemize}
    \item Frontend: menerima input dan menampilkan hasil
    \item Backend: melakukan parsing HTML dan traversal DOM
\end{itemize}
